\documentclass[12pt,a4paper]{article}
\usepackage[margin=25mm, headheight=15pt, headsep=10pt]{geometry}
\usepackage{fontspec}
\usepackage[table]{xcolor} % Note: table option is required for rowcolors
\usepackage{titlesec}
\usepackage{tocloft}
\usepackage{fancyhdr}
\usepackage{booktabs}
\usepackage{tcolorbox}
\tcbuselibrary{skins, breakable}
\usepackage[colorlinks=true, linkcolor=emerald, urlcolor=emerald, bookmarks=true]{hyperref}

\definecolor{emerald}{RGB}{0,102,51}
\definecolor{darkred}{RGB}{139,0,0}
\definecolor{lightgray}{RGB}{245,245,245}

\usepackage[nil]{babel}
\babelprovide[import, main]{bengali}
\babelprovide[import]{arabic}
\babelfont{rm}[Renderer=HarfBuzz,Scale=1.0]{Noto Serif Bengali}
\babelfont{sf}[Renderer=HarfBuzz]{Noto Sans Bengali}
\babelfont{tt}[Renderer=HarfBuzz]{Noto Sans Bengali}
\babelfont[arabic]{rm}[Renderer=HarfBuzz,Scale=1.1]{Noto Naskh Arabic}

% Section styling
\titleformat{\section}{\Large\bfseries\color{emerald}}{\thesection.}{0.5em}{}
\titleformat{\subsection}{\large\bfseries\color{emerald!85!black}}{\thesubsection.}{0.5em}{}
\titleformat{\subsubsection}{\normalsize\bfseries\color{emerald!70!black}}{\thesubsubsection.}{0.5em}{}
\titlespacing*{\section}{0pt}{18pt}{8pt}

% Fancy Header/Footer setup
\pagestyle{fancy}
\fancyhf{} % clear all header and footer fields
\fancyhead[L]{\color{emerald}\small\textbf{\nouppercase{\leftmark}}}
\fancyfoot[L]{\scriptsize\color{black!50}{\lat{AI}}-সংকলিত, আলেম-যাচাইকৃত নয়}
\fancyfoot[C]{\color{emerald!70!black}\thepage}
\renewcommand{\headrulewidth}{0.4pt}
\renewcommand{\headrule}{\hbox to\headwidth{\color{emerald}\leaders\hrule height \headrulewidth\hfill}}
\renewcommand{\footrulewidth}{0pt}

% Custom boxes
% 1. Ayat/Hadith Box
\newtcolorbox{ayatbox}[1][]{
    enhanced, breakable, 
    colback=emerald!5, 
    colframe=emerald, 
    boxrule=0.5pt, 
    arc=3pt, 
    left=10pt, right=10pt, top=10pt, bottom=10pt,
    #1
}

% 2. Quote/Translation Box
\newtcolorbox{quotebox}[1][]{
    enhanced, breakable,
    frame hidden,
    colback=lightgray,
    borderline west={3pt}{0pt}{emerald},
    left=15pt, right=10pt, top=10pt, bottom=10pt,
    sharp corners,
    #1
}

% 3. AI-generated notice box
\newtcolorbox{warnbox}[1][]{
    enhanced, breakable,
    colback=red!4,
    colframe=darkred,
    boxrule=0.8pt,
    arc=3pt,
    left=10pt, right=10pt, top=10pt, bottom=10pt,
    #1
}

% Commands
\newcommand{\arab}[1]{\foreignlanguage{arabic}{#1}}
\newcommand{\kib}[1]{\textcolor{emerald}{\textbf{#1}}}

% Latin (English) fallback: Noto Serif Bengali has NO Latin glyphs
\newfontfamily\latfont[Renderer=HarfBuzz]{Noto Serif}
\newcommand{\lat}[1]{{\latfont #1}}

\setlength{\parskip}{5pt}
\setlength{\parindent}{0pt}

% Fix for missing bullet in Noto Serif Bengali
\renewcommand{\labelitemi}{$\bullet$}



\begin{document}
\begin{center}{\Large\bfseries\color{emerald} ঘটনা ৪ — হুজাইফা (রাঃ): নবীজির পেছনে দীর্ঘ রাতের নামাজ}\end{center}
\vspace{1em}
\section*{ঘটনা ৪ — হুজাইফা (রাঃ): নবীজির পেছনে দীর্ঘ রাতের নামাজ}

\begin{warnbox}
 \textbf{গুরুত্বপূর্ণ নোটিশ (\lat{Important} \lat{Notice}):} এই কন্টেন্টটি \lat{AI} (কৃত্রিম বুদ্ধিমত্তা) দিয়ে প্রস্তুত ও সংকলিত। \lat{AI} কঠোর সূত্র-মান নীতি মেনে শুধু নির্ভরযোগ্য উৎস থেকে তথ্য সংগ্রহ করেছে — কেবল সহীহ/হাসান হাদিস ও সহীহ সনদের আসার রাখা হয়েছে, দুর্বল সনদ বাদ দেওয়া হয়েছে। তবে এটি কোনো আলেম বা মুহাদ্দিস কর্তৃক যাচাইকৃত নয়।
\end{warnbox}

\begin{quote}
\textbf{সম্পর্কিত ফাইল:} পার্ট ৫ঘ (তাহাজ্জুদ-নফল), পার্ট ৪ (পদ্ধতি)।
\end{quote}

\medskip\hrule\medskip

\subsection*{১. সূত্র কার্ড (\lat{Source} \lat{Card})}

\begin{table}[h]\centering\small
\begin{tabular}{ll}
বিষয় & বিবরণ \\
\hline
\textbf{ঘটনা} & হুজাইফা (রাঃ) রাতের নামাজে নবীজির (সা.) পেছনে দীর্ঘতম কিরাআত-রুকু-সিজদা প্রত্যক্ষ করেন \\
\textbf{রাবি} & হুজাইফা ইবনুল ইয়ামান (রাঃ) \\
\textbf{গ্রন্থ ও নম্বর} & সহীহ মুসলিম ৭৭২; সুনানে নাসাঈ ১৬৬৪ \\
\textbf{সনদের মান} & সহীহ \\
\hline
\end{tabular}
\end{table}

\medskip\hrule\medskip

\subsection*{২. পটভূমি (\lat{Background})}

\textbf{কে:} নবীজি মুহাম্মাদ (সা.); হুজাইফা ইবনুল ইয়ামান (রাঃ) — নবীজির সিরের (গোপন বিষয়) জ্ঞানী সাহাবি; কুরআন-মুখস্থ, গভীর ঈমানদার।
\textbf{কখন:} মদিনা যুগের এক রাত — বিশেষত রমজান-উত্তর/সাধারণ রাত বলা হয়; বর্ণনায় সুনির্দিষ্ট তারিখ নেই।
\textbf{কোথায়:} মদিনা — মসজিদে নববি অথবা একটি খেজুর-শাখার ঘরে।
\textbf{কেন:} হুজাইফা (রাঃ) নবীজির (সা.) রাতের নামাজের প্রকৃত পদ্ধতি — দীর্ঘ কিরাআত, রুকু-সিজদার মাত্রা — প্রত্যক্ষ করবার জন্য।

\medskip\hrule\medskip

\subsection*{৩. পূর্ণ ঘটনার বর্ণনা (\lat{The} \lat{Full} \lat{Event})}

হুজাইফা (রাঃ) বলেন:

\begin{quote}
\lat{“}এক রাতে আমি নবীজির (সা.) সাথে নামাজ পড়লাম। তিনি সূরা আল-বাকারাহ শুরু করলেন। আমি ভাবলাম — ১০০ আয়াত পেরোলেই রুকু করবেন — কিন্তু তিনি চালিয়ে গেলেন। আমি ভাবলাম — এক রাকাতে পুরো সূরা পড়বেন — কিন্তু তাও অতিক্রম করলেন। আমি ভাবলাম — সূরা শেষ করে রুকু করবেন — তিনি আবার সূরা আন-নিসা শুরু করলেন এবং সম্পূর্ণ পড়লেন; তারপর সূরা আল-ইমরান — পুরোটা, ধীরে-ধীরে (তিলাওয়াত)। এমনকি – যখন রহমতের আয়াত পের হতেন, তিনি রহমত চাইতেন; শাস্তির আয়াত পের হতেন, আউজুবিল্লাহ পড়তেন (পাহারা চাইতেন); তাসবিহের আয়াত পের হতেন, তাসবিহ পড়তেন।\lat{”}
\end{quote}

তারপর নবীজি (সা.) রুকু করেন ও \textbf{\lat{“}সুবহানা রাব্বিয়াল আযিম\lat{”}} পড়েন — \textbf{রুকু প্রায় দাঁড়ানোর সমান দীর্ঘ} ছিল। মাথা উঠিয়ে \textbf{\lat{“}সামিআল্লাহু লিমান হামিদাহ, রাব্বানা লাকাল হামদ\lat{”}} — বল\arab{اشر}ক আবার \textbf{দাঁড়িয়ে থাকলেন দাঁড়ানোর সমান} (কিয়াম-পরবর্তী স্থিতি)। তারপর সিজদা ও \textbf{\lat{“}সুবহানা রাব্বিয়াল আলা\lat{”}} — \textbf{সিজদাও প্রায় দাঁড়ানোর সমান} — ছিল। হুজাইফা (রাঃ) বিস্ময়ের সঙ্গে বললেন:

\begin{quote}
\lat{“}তিনি এক রাকাতে আল-বাকারাহ, আন-নিসা ও আল-ইমরান পড়েছেন; রুকু ও সিজদা — প্রত্যেকটি প্রায় কিয়ামের সমান। এত দীর্ঘ ও গভীর নামাজ আমি আর দেখিনি।\lat{”}
\end{quote}

বর্ণনায় আরও আছে — নবীজি (সা.) \textbf{রাতের নামাজে চার রাকাত পর্যন্ত} পড়েছেন — তাহাজ্জুদের পদ্ধতি; এবং রমজানের রাতে খেজুর-ঘরে ১১ রাকাতের (৮+৩ বিতর) আমলও হুজাইফার (রাঃ) পর্যবেক্ষণে আছে (নাসাঈ ১৬৬৪-এর আলোকে)।

\medskip\hrule\medskip

\subsection*{৪. শব্দের ব্যাখ্যা (\lat{Key} \lat{Words})}

\begin{itemize}
\item \textbf{তারতীল:} ধীরে, সুস্পষ্ট, অর্থ-সহ উপলব্ধিসহ তিলাওয়াত।
\item \textbf{তুমাঈনাহ:} রুকন-আন্দোলনের মধ্যে স্থিরতা/পূর্ণতা — নামাজের আবশ্যকতা (সব মাযহাব)।
\item \textbf{আউজুবিল্লাহ:} \lat{“}আল্লাহর কাছে পাহারা চাই\lat{”} — শাস্তির আয়াত পড়লে।
\item \textbf{সামিআল্লাহু লিমান হামিদাহ / রাব্বানা লাকাল হামদ:} রুকু থেকে ওঠার জিকির — \lat{“}আল্লাহ শোনেন-স্বীকার করেন প্রশংসাকারীর; হে আমাদের রব, তোমারই প্রশংসা।\lat{”}
\end{itemize}
\medskip\hrule\medskip

\subsection*{৫. সংশ্লিষ্ট হাদিস ও আয়াত}

\begin{quote}
\textbf{\lat{“}অতঃপর কুরআন পড়ো নিয়ম অনুযায়ী (তারতীল)।\lat{”}} — সূরা আল-মুজ্জাম্মিল ৭৩:৪ (বাংলা অর্থ)।
\end{quote}

\begin{quote}
\textbf{\lat{“}যারা রাতে নামাজে দাঁড়ায় এবং সিজদা করে — আল্লাহ তাদের প্রতি রহমত বর্ষণ করেন।\lat{”}} — সূরা আল-ফুরকান ২৫:৬৪-এর অর্থ-ভিত্তি (বাংলা অনুবাদ)।
\end{quote}

\begin{quote}
নবীজি (সা.) বলেন: \textbf{\lat{“}যে চোখ রাতের অন্ধকারে (নামাজের জন্য) জাগে — যেমন দুটি চোখ, যা আর কখনো দেখবে না — অর্থাৎ জান্নাতের প্রদীপ।\lat{”}} — (সুনানে তিরমিজি ২৩৩-এর আলোকে)।
\end{quote}

\medskip\hrule\medskip

\subsection*{৬. রেফারেন্স}

\begin{itemize}
\item সহীহ মুসলিম — ৭৭২
\item সুনানে নাসাঈ — ১৬৬৪
\item সুনানে তিরমিজি — ২৩৩ (রাতের নামাজের ফজিলত)
\item কুরআন: আল-মুজ্জাম্মিল ৭৩:৪; আল-ফুরকান ২৫:৬৪
\end{itemize}

\end{document}
